\documentclass[12pt]{article}

\usepackage{amsmath, amssymb}
\usepackage{geometry}
\usepackage[english]{babel}

\geometry{a4paper, margin=25mm}

% タイトルから不要な「The」や「A」の重複を整理し、格調高く修正
\title{The 0th Overview of the Wave-Based Interpretation of Spatial Structure}
\author{Night Sleepless \\ Independent Researcher}
\date{\today}

\begin{document}

\maketitle

\begin{abstract}
This paper introduces the Wave-Based Interpretation of Spatial Structure, a minimal framework in which space itself is treated as a multi-scale wave field possessing intrinsic periodicity and phase structure. 
In this interpretation, galactic rotation curves, disk formation, coherent spin alignment, and the non-destructive nature of galactic encounters arise naturally from the superposition of spatial status. 
Stars and stellar systems do not require additional mass components to remain bound; instead, they follow the stable phase patterns of the underlying wave field. 
The apparent mass discrepancy commonly attributed to dark matter is reinterpreted as an observational omission of higher-order spatial status. The model is presented in a dialogue-based summary format to preserve the reasoning process and highlight the conceptual simplicity of treating gravitational behavior as emergent from wave interference rather than invisible matter. 
This approach suggests that resolving the structure of space itself may provide a unified explanation for several outstanding problems in astrophysics and cosmology.
At the core of the model lies the relation \( m \propto 1/\tau_{\text{space}} \) which links mass to the characteristic periodicity of space itself.
\end{abstract}

\section{Introduction}

% 「form」を「from」に、「th」を「the」に修正
This document summarizes a framework arising from a dialogue-based exploration of the structure of space. Rather than introducing new elementary particles or hidden dimensions to reconcile astronomical anomalies, this approach examines whether longstanding discrepancies in galactic dynamics can be resolved by analyzing the intrinsic properties of the spatial medium itself. 

Historically, physical models have treated space as a passive, non-periodic background in which discrete matter particles interact. When large-scale observations—such as galactic rotation velocities—deviate from classical predictions based on visible matter, conventional physics introduces non-baryonic dark matter to preserve the underlying gravitational equations. While mathematically functional within specific bounds, this method introduces persistent conceptual difficulties regarding the physical detection and lifecycle of such unobservable mass components.

The framework discussed here suggests that the missing mass is an artifact of treating space as an invariant backdrop. By shifting the focus to the wave characteristics of space, observed gravitational behaviors are reinterpreted as phase-driven phenomena emergent from spatial status.

\section{The Core Relation: Mass and Spatial Periodicity}

The foundational premise of this framework is that observable mass \( m \) is not an inherent, isolated property of a particle, but rather a localized manifestation of spatial resistance. This relationship is governed by the characteristic periodicity of space, \( \tau_{\text{space}} \), expressed as:
\[
m \propto \frac{1}{\tau_{\text{space}}}
\]
In this model, what is conventionally measured as mass corresponds to the phase reaction or restorative tension generated when external momentum is applied to a periodic spatial structure. A smaller spatial period \( \tau_{\text{space}} \) implies a higher structural density or tighter tension within the spatial field, yielding a greater macroscopic inertial or gravitational response (mass). 

This perspective avoids the necessity of active field excitation by a particle to generate mass. Instead, mass is an inherent, real-time reactive consequence—an action-reaction counter-balance—of the spatial wave field maintaining its equilibrium configuration against localized kinetic stress.

\section{Galactic Dynamics Without Dark Matter}

Applying the relation \( m \propto 1/\tau_{\text{space}} \) to galactic scales provides an alternative explanation for flat rotation curves without invoking halos of dark matter. 

Galaxies are treated as coherent macroscopic wave systems where the underlying spatial medium exhibits structured phase architecture. The distribution of stars and gas is governed by the stable nodes and troughs of this multi-scale spatial wave field. Stellar systems do not require gravitational binding from invisible matter; they follow the localized phase trajectories established by the spatial geometry itself.

Furthermore, this framework naturally accounts for the observed coherence in stellar disk alignment and the non-destructive nature of galactic encounters. When galaxies interact or pass through one another, they are not colliding as massive particle aggregates, but rather as overlapping wave states. The systems transition through superposition, retaining their structural memory and disk integrity because their organizational patterns are anchored in the persistent periodicity of the spatial medium.

\section{Dialogue-Based Development Summary}

The concepts outlined in this overview were developed through a systematic, dialogue-based analytical process designed to isolate core physical principles from historical mathematical assumptions. The reasoning process progressed through three distinct phases:

\begin{enumerate}
    \item \textbf{Deconstruction of Particle-Centric Assumptions:} Analyzing the structural vulnerabilities of dark matter models, specifically the logical and energetic paradoxes introduced by treating gravitational anomalies as discrete, hidden particles.
    \item \textbf{Inversion of Background and Foreground:} Shifting the operational framework from matter acting inside an empty space to space acting as a dynamic, periodic medium whose geometric states dictate material behavior.
    \item \textbf{Axiomatic Formulation:} Condensing the resulting insights into a single, scalable relation (\( m \propto 1/\tau_{\text{space}} \)) capable of bridging local inertia with cosmic structures without requiring parametric fine-tuning.
\end{enumerate}

This collaborative, dialogue-driven distillation process ensured that each conceptual layer remained strictly bounded by classical conservation laws while systematically removing non-essential variables.

\section{Discussion and Future Inferences}

The formal adoption of spatial periodicity opens immediate pathways for resolving higher-order cosmological discrepancies. Subsequent developments of this model will expand upon the specific geometric topology of the spatial wave field, exploring how specific periodic configurations—such as uniform isotropic lattices—constrain particle lifetimes and define the limits of energy conservation during localized high-energy impacts.

Additionally, this approach provides a rigid mechanism for evaluating cosmic expansion and redshift. If space possesses an intrinsic base oscillation, the systematic shifting of distant light can be analyzed as a propagation characteristic—a historical record of phase interaction over extreme distances—rather than a physical stretching of the background canvas. Future iterations will formalize these kinematics to explicitly map out the behavior of higher-generation states.

By organizing the framework into succinct, geocentrically stable propositions, the model remains accessible to immediate empirical validation through existing galactic velocity datasets, because stars follow locally stable phase configurations rather than relying on direct mass-based binding.

More broadly, the structure suggests that many cosmological phenomena may be better understood by analyzing the structure and periodicity of the space itself. Treating space as a wave field provides a conceptual foundation that unifies galactic dynamics, gravitational behavior, and large-scale coherence under a single principle.

\section{Conclusion}

The space periodicity proposes that gravitational behavior and galactic structure arise from the intrinsic periodicity and phase architecture of space itself. By treating space as a multi-scale wave field, the structure provides a unified explanation for rotation curves, disk coherence, and the non-destructive nature of galactic encounters without invoking additional invisible mass components. The core relation \( m \propto 1/\tau_{\text{space}} \) serves as the foundation that links observable mass to the underlying spatial oscillation.

% 「qality」を「quality」に修正
This perspective suggests that many longstanding astrophysical discrepancies may originate not from missing matter but from an incomplete characterization of spatial quality. Resolving the structure of space therefore becomes central to understanding gravitational dynamics and large-scale cosmic organization. The dialogue-based development summary preserved in this work highlights the conceptual simplicity of the model and provides a pathway for future refinement and observational evaluationｐ.

\end{document}
